% =============================================================================
% Meeting handout: geometric spring correction (theory + code changes)
%
% Compile:
%   cd papers/technical && latexmk -pdf -output-directory=../build -aux-directory=../build spring_force_geometric_fix_note.tex
% =============================================================================
%
% Purpose: very brief, professor-facing note explaining
% (1) what was wrong with the prior displacement-based direction,
% (2) the corrected geometric formulation, and
% (3) concrete code changes needed to match the corrected formulation.
%
% This file is intentionally short.
% =============================================================================

\documentclass[11pt,a4paper]{article}
\usepackage[utf8]{inputenc}
\usepackage{amsmath,amssymb,amsfonts}
\usepackage[margin=1in]{geometry}
\usepackage{hyperref}
\usepackage{booktabs}
\usepackage{microtype}
\usepackage{parskip}
\usepackage{enumitem}
\usepackage{tcolorbox}
\tcbuselibrary{breakable}

\setlist{itemsep=0.25em, topsep=0.25em}

\newtcolorbox{takeaway}{colback=gray!7,colframe=gray!55,boxrule=0.6pt,arc=2pt,left=7pt,right=7pt,top=6pt,bottom=6pt,breakable}

\title{\textbf{Geometric spring-force correction (meeting note)}}
\author{Daniel Miles}
\date{\today}

\begin{document}
\maketitle

\section{What the feedback means (one sentence)}
A physical spring force should act along the \emph{current line connecting its endpoints} (the deformed chord), but our previous write-up/code used the \emph{relative displacement} vector to set the force direction, which can point in a different direction under shear-like motion.

\section{The key distinction}
Let a bond connect nodes $A$ and $B$.

\textbf{Reference (undeformed) geometry:}
\[
\mathbf{X}_A^{eq},\ \mathbf{X}_B^{eq},\qquad \mathbf{R}_0 = \mathbf{X}_B^{eq}-\mathbf{X}_A^{eq},\qquad L_0 = \|\mathbf{R}_0\|.
\]

\textbf{Displacements stored by the solver:}
\[
\mathbf{u}_A,\mathbf{u}_B,\qquad \mathbf{X}_A=\mathbf{X}_A^{eq}+\mathbf{u}_A,\quad \mathbf{X}_B=\mathbf{X}_B^{eq}+\mathbf{u}_B.
\]

\textbf{Current (deformed) chord in space:}
\begin{equation}
\mathbf{R}=\mathbf{X}_B-\mathbf{X}_A=\mathbf{R}_0 + (\mathbf{u}_B-\mathbf{u}_A),\qquad L=\|\mathbf{R}\|,\qquad \hat{\mathbf{R}}=\mathbf{R}/L.
\label{eq:Rdef}
\end{equation}

\textbf{What we previously used:}
\[
\mathbf{r}_{old}=\mathbf{u}_B-\mathbf{u}_A,\qquad \hat{\mathbf{r}}_{old}=\mathbf{r}_{old}/\|\mathbf{r}_{old}\|.
\]

\begin{takeaway}
\textbf{Why the professor's counterexample works.} If the reference bond is horizontal ($\mathbf{R}_0=\hat{\mathbf{i}}$) but only $B$ moves up ($\mathbf{u}_B=\hat{\mathbf{j}}$), then $\mathbf{R}=\hat{\mathbf{i}}+\hat{\mathbf{j}}$ (45$^\circ$), while $\mathbf{r}_{old}=\hat{\mathbf{j}}$ (90$^\circ$). A geometric spring must push/pull along $\hat{\mathbf{R}}$, not along $\hat{\mathbf{r}}_{old}$.
\end{takeaway}

\section{Corrected spring force (geometric)}
Define the scalar extension/compression
\begin{equation}
s = L - L_0.
\label{eq:sdef}
\end{equation}

One consistent exponential analogue of Hooke's law is
\begin{equation}
\mathbf{F}_{A\leftarrow B} = k\left(e^{\alpha s}-1\right)\hat{\mathbf{R}}.
\label{eq:Fgeo}
\end{equation}

Notes:
\begin{itemize}
  \item For small $|s|$, $e^{\alpha s}-1\approx \alpha s$, so $F\approx (k\alpha)s$.
  \item If $s<0$ (compression), then $e^{\alpha s}-1<0$ and the force reverses direction, as expected.
  \item The direction is always along the current chord $\hat{\mathbf{R}}$ (geometrically consistent).
\end{itemize}

\section{What must change in the write-up}
In \texttt{papers/technical/spring\_force\_11x11.pdf}, the spring section should be revised so that:
\begin{itemize}
  \item the vector used for direction is $\mathbf{R}$ in Eq.~\eqref{eq:Rdef}, not $\mathbf{u}_B-\mathbf{u}_A$;
  \item the scalar ``stretch'' entering the exponential is $s=L-L_0$ in Eq.~\eqref{eq:sdef} (or an equivalent measure), not $d=\|\mathbf{u}_B-\mathbf{u}_A\|$;
  \item nearest-neighbor bonds use $L_0=\Delta$ and diagonal bonds use $L_0=\Delta\sqrt{2}$, where $\Delta$ is the grid spacing.
\end{itemize}

\section{What must change in the code (minimal, concrete)}
Current behavior (in \texttt{src/lattice\_simulation\_11x11.jl}) computes the spring direction from \texttt{pos2-pos1}, where \texttt{pos} is the displacement state. To match Eq.~\eqref{eq:Fgeo} we must include the reference chord $\mathbf{R}_0$ for each neighbor link.

\textbf{Minimal implementation strategy:}
\begin{itemize}
  \item In the neighbor loop (where we know whether a link is left/right/up/down/diagonal), define $\mathbf{R}_0$ from the neighbor offset:
  \[
  \mathbf{R}_0 =
  \begin{cases}
  [\pm\Delta,0]^T & \text{horizontal},\\
  [0,\pm\Delta]^T & \text{vertical},\\
  [\pm\Delta,\pm\Delta]^T & \text{diagonal},
  \end{cases}
  \qquad L_0=\|\mathbf{R}_0\|.
  \]
  \item Compute $\mathbf{R}=\mathbf{R}_0+(\mathbf{u}_B-\mathbf{u}_A)$, then $s=\|\mathbf{R}\|-L_0$, then apply Eq.~\eqref{eq:Fgeo}.
  \item Replace the current call to \texttt{spring\_force\_2d(pos\_k, pos\_neighbor, k, alpha)} with a new routine such as \texttt{spring\_force\_2d\_geometric(uA,uB,R0,k,alpha)}.
\end{itemize}

\textbf{Also update:}
\begin{itemize}
  \item the spring potential-energy function to use $s$ and $L_0$ if you want energy checks consistent with the new dynamics;
  \item any text in the manuscript that claims ``force acts along the line connecting the masses'' (this becomes true again once Eq.~\eqref{eq:Fgeo} is used).
\end{itemize}

\section{Suggested discussion points for the meeting (what to show)}
\begin{itemize}
  \item The one-line correction: ``use the deformed chord $\mathbf{R}=\mathbf{R}_0+(\mathbf{u}_B-\mathbf{u}_A)$ for direction, and $s=\|\mathbf{R}\|-L_0$ for the scalar input.''
  \item The professor's $(4,4)$--$(5,4)$ example and the 45$^\circ$ chord.
  \item A clear statement of what changes in simulation behavior: shear no longer produces a force direction unrelated to the actual spring axis.
  \item Clarify whether the model should allow compression ($s<0$) with the same exponential law or whether you prefer a tension-only/contact-like behavior; Eq.~\eqref{eq:Fgeo} allows both.
\end{itemize}

\end{document}

